% !PRINTED 718
\weberpage{718}

Demnach ergiebt sich aus \S.\,194, (12) mit Rücksicht auf die trigonometrische Formel
\[
\sin\frac{3\pi}{8} = \cos\frac{\pi}{8}
\]
\[
X_1 = -\frac{1}{\sqrt{2}} \log \operatorname{tg}\frac{\pi}{8}, \quad X_2 = \frac{-\pi}{4}, \quad X_3 = \frac{-\pi}{2\sqrt{2}}.
\]
Nun ist
\[
\tau = \frac{r-1}{r^2(r+1)} = \operatorname{tg}\frac{\pi}{8} = \sqrt{2} - 1
\]
eine Einheit des Körpers $H_8$, der hier nichts Anderes ist, als der aus $\sqrt{2}$ entspringende quadratische Körper, und die Einheiten darin erhält man also aus den Lösungen der \gesperrt{Pell}'schen Gleichung
\[
u^2 - 2v^2 = \pm 1
\]
in der Form $u + v\sqrt{2}$.

Die kleinste positive Lösung dieser \gesperrt{Pell}'schen Gleichung ist aber $u = 1$, $v = 1$, und folglich findet man nach Bd.\,I, \S.\,128 alle Einheiten in $H_8$ in der Form $\pm(1+\sqrt{2})^a$, worin $a$ ein positiver oder negativer ganzzahliger Exponent ist; $\tau$ ergiebt sich hieraus für $a = -1$, und daher sind alle Einheiten auch in der Form
\[
\pm \tau^a
\]
enthalten. Es ist also $\tau$ eine \gesperrt{fundamentale Einheit}, und die in den Formeln des vorigen Paragraphen vorkommende Determinante $E$ [\S.\,195, (12)] wird hier (da $\tau$ ein echter Bruch ist):
\[
E = -\log \tau.
\]
Man erhält also aus \S.\,195, (15), (18):
\[
h_1 = 1, \quad k = 1, \quad h = 1.
\]
Der Körper $\Omega_8$ ist also ein \gesperrt{einclassiger}, d.\,h. ein solcher, in dem die Zerlegung der ganzen Zahlen in Primfactoren nach denselben Gesetzen, wie im Körper der rationalen Zahlen, möglich ist.

Es hat sich dabei zugleich ergeben, dass alle Einheiten im Körper $H_8$ in der Form $\pm \tau^a$ darstellbar sind. Die conjugirten Werthe der Einheit $\tau$ sind aber
\[
\tau_1 = \operatorname{tg}\frac{\pi}{8}, \quad \tau_3 = -\operatorname{tg}\frac{3\pi}{8} = -\tau_1^{-1},
\]
und haben also verschiedene Zeichen. Wir schliessen daraus für diesen Fall auf den Satz:
%CONT>>
